% ============================================================
%  前言 双语导读 —— Preface（含 Foreword 提要）
%  形式：中文概述 + 关键原句短引用 + 解读 + 实践清单
% ============================================================

\chapter*{Preface\quad 前言导读}
\addcontentsline{toc}{chapter}{Preface 前言（导读）}
\markboth{Preface}{前言}

\noindent{\small 前言不长，但把全书的「世界观」交代清楚了：这是一本关于\emph{做事}
（doing）的书，目标是让你成为更好的程序员。下面按原书的小标题逐段导读。}

\vspace{0.5em}

\begin{bitext}
\src{This book will help you become a better programmer.}
\tgt{这本书会帮助你成为一名更好的程序员。}
\end{bitext}

\commentary{作者一上来就定调：无论你是独立开发者、大团队的一员，还是同时服务多个
客户的顾问，这本书都是写给\emph{个体}的，帮你把工作做得更好。它不空谈理论，而是聚
焦实践。作者还做了个词源小考：pragmatic 源自拉丁语 pragmaticus（「精于事务」），
再往上溯自希腊语，意为「去做」。所以这是一本关于「去做」的书。}

\begin{bitext}
\src{Programming is a craft.}
\tgt{编程是一门手艺。}
\end{bitext}

\commentary{作者把程序员描述成四种角色的混合体：一部分倾听者、一部分顾问、一部分
翻译、一部分「独裁者」。你要在项目时钟的无情滴答声中，捕捉飘忽的需求、写好文档、
把工程做扎实——「你每天都在创造小小的奇迹」。紧接着作者泼了盆冷水：工具厂商、方法
论大师人人都号称自己的东西是万能解药，但
「There are no easy answers.」没有银弹，也不存在普适的「最佳方案」，只有在特定
情境下更合适的方案。}

% ------------------------------------------------------------
\sechead{Why ``Pragmatic''?}{为什么叫「务实」}

\begin{ideabox}
务实（pragmatism）意味着：不要死守某一种技术，而要有足够广的背景与经验，好在具体
情境下挑出合适的解。理论（计算机科学的基本原理）与实践（各式各样的真实项目）合起来
才让你强大。你的做法要随环境和情境不断调整，并在工作推进中\emph{持续}做这种判断。
\end{ideabox}

\commentary{这一段是全书价值观的浓缩。作者反对「一招鲜」，也反对纯粹的学院派——强调
经验与判断力。一句话概括：「Pragmatic Programmers get the job done, and do it well.」
把事情做成，并且做好。}

% ------------------------------------------------------------
\sechead{Who Should Read This Book?}{谁该读这本书}

\begin{ideabox}
写给想变得\emph{更高效、更有产出}的程序员。如果你觉得自己没发挥出应有的潜力，或者
看着同事用工具比你更高效，或者你当前的工作用的是老技术、想知道新思路怎么用上——这本
书就是给你的。作者很坦诚：他们并不掌握全部（甚至不是大部分）答案，想法也不都适用于
所有情况；但只要你跟着走，经验会快速增长、效率会提高、对开发全流程的理解会更透彻。
\end{ideabox}

% ------------------------------------------------------------
\sechead{What Makes a Pragmatic Programmer?}{什么造就了务实程序员}

\begin{ideabox}
每位开发者都是独一无二的，会形成自己的个人环境。但务实程序员共享以下特质：
\textbf{早期采用者/快速适应者}（对新事物有直觉，乐于尝试，上手快）、
\textbf{好奇}（爱发问，像囤积癖一样收藏小事实）、
\textbf{批判性思考者}（不轻信「一直都是这么做的」，对厂商承诺保持警惕）、
\textbf{务实}（努力理解问题的本质，因而对难度和耗时有良好判断）、
\textbf{多面手}（尽量熟悉广泛的技术，能随时转向新领域）。
最后是两条最基本的、所有务实程序员都具备的特质，作者直接把它们列为 Tip 1 和 Tip 2。
\end{ideabox}

\begin{bitext}
\src{\tip{1}{Care About Your Craft}{We feel that there is no point in
developing software unless you care about doing it well.}}
\tgt{\tip{1}{关心你的手艺}{我们认为，除非你在意把软件做好，否则开发软件毫无意义。}}
\end{bitext}

\begin{bitext}
\src{\tip{2}{Think! About Your Work}{We're challenging you to think about what
you're doing while you're doing it. Never run on auto-pilot.}}
\tgt{\tip{2}{思考！关于你的工作}{我们要求你在做事的同时思考你在做的事。永远不要开启
「自动驾驶」。}}
\end{bitext}

\commentary{Tip 2 是全书的方法论内核：\emph{实时地}批判、审视自己的工作，而不是事后
做一次性的复盘。作者用 IBM 那句老口号「THINK!」作结，称其为务实程序员的箴言。他们还
很诚实地承认这套做法\emph{很累}——要占用你本就紧张的宝贵时间；回报则是更投入地从事
你热爱的工作、对越来越多领域有掌控感，以及持续进步的愉悦。长期看，你和团队会更高效、
代码更易维护、开会更少。}

% ------------------------------------------------------------
\sechead{Individual Pragmatists, Large Teams}{个体的务实主义者，庞大的团队}

\begin{ideabox}
有人主张大团队/复杂项目里没有个人发挥的空间，说「软件构建是一门工程学科，个体自作主张
就会崩盘」。作者不同意：软件构建\emph{理应}是工程学科，但这并不排斥个人的手艺。他们用
中世纪大教堂作比——成千上万人年的投入、跨越数十年，木匠、石匠、雕刻匠、玻璃匠都是
手艺人，他们在工程要求之上做出了超越纯机械层面的整体。支撑项目走下去的，正是他们对
自身贡献的信念。
\end{ideabox}

\begin{bitext}
\src{We who cut mere stones must always be envisioning cathedrals.
\upshape\small\hfill---Quarry worker's creed}
\tgt{切割石头的人，心中始终要想着大教堂。\hfill---采石工人的信条}
\end{bitext}

\commentary{作者的判断是：在一百年的尺度上，今天我们的「工程」也许会被后人看作像中世纪
工匠的方法一样古老，而我们的\emph{手艺}仍将被尊重。换言之，工程规范与个人技艺不是对立
的，后者才是长期价值所在。}

% ------------------------------------------------------------
\sechead{It's a Continuous Process}{这是一个持续的过程}

\begin{ideabox}
访客问伊顿公学的园丁怎么把草坪养得那么完美，园丁答：「每天早上拂去露水，隔天割一次草，
每周滚压一次——这样坚持 500 年，你也会有一片好草坪。」优秀的草坪需要少量而\emph{持续}
的照顾，优秀的程序员也一样。作者引入日语词 \emph{kaizen}（改善）：持续做大量小的改进，
这是日本制造业质量与生产率大幅提升的主因之一。它同样适用于个人：每天打磨已有技能、给
自己的工具箱添点新家伙。
\end{ideabox}

\commentary{这节和第 1 章的「知识资产」遥相呼应。园丁的段子点出关键：不是靠某次爆发，
而是靠\emph{微量且不停}的投入。作者还给了个乐观的对比——不像伊顿草坪要等 500 年，你的
努力几天内就能看到效果，几年后会惊讶于自己的成长。}

% ------------------------------------------------------------
\sechead{How the Book Is Organized}{本书如何组织}

\begin{ideabox}
全书由一系列\emph{短小、自足}的小节构成，每节讨论一个主题，彼此有大量交叉引用。可以
任意顺序阅读，不必从头读到尾。文中会出现标为 Tip nn 的提示框；作者说这些 tip「有自己的
生命」，他们自己每天都在照做。附录 A 汇集资源（参考书目、网址、推荐期刊/书籍/专业组织），
附录 B 是练习答案。书中还穿插 exercise（答案较直接）与 challenge（更开放），后者可用于
小组讨论或课程作业。
\end{ideabox}

% ------------------------------------------------------------
\sechead{What's in a Name?}{名词之辨}

\begin{ideabox}
书里会冒出各种术语，有的是被「技术化」的普通英文词，有的是计算机科学家生造的怪词。作者
尽量在首次出现时给出定义或提示，但难免有漏网的。作者给读者的建议很实在：遇到没见过的词
别跳过，花点时间查一查（上网或翻教材），查不到就发邮件来吐槽，好让他们在下一版补上定义。
\end{ideabox}

\commentary{这节还有个有趣的点：作者「反击」了计算机科学家——他们有时会\emph{故意}不用
现成的术语，因为那些术语通常被绑死在某个特定问题域或开发阶段，而本书主张大多数技术在
代码、设计、文档、团队组织之间是通用的。现成术语一放到更广的语境就产生歧义，于是作者
干脆自造新词。（Humpty Dumpty 那句「一个词的意思就是我选定它是什么意思」是这节的题眼。）}

% ------------------------------------------------------------
\sechead{Source Code, Feedback \& Acknowledgments}{源码、反馈与致谢}

\begin{ideabox}
书中大部分代码取自可编译的源文件，可在 \texttt{www.pragmaticprogrammer.com} 下载，
站上还有作者觉得有用的资源链接、勘误与更新。欢迎读者反馈（评论、建议、正文错误、
示例问题），邮箱 \texttt{ppbook@pragmaticprogrammer.com}。致谢部分作者感谢了
Addison-Wesley 的编辑与制作团队，以及一长串审阅者——并打趣说，没有这些细致的意见，
这本书会更难读、更不准确、\emph{还要长一倍}。
\end{ideabox}

\vspace{0.6em}
\begin{center}
\small\itshape\color{ppgray}
前言导读完。全书 70 条 TIP 的完整标题对照见附录 A。
\end{center}
